We presented scRareRefine, a post-hoc inductive rescue framework for rare cell type annotation in semi-supervised single-cell data. The core contribution is the separability ratio $S = \delta_r / \rho_r$, a scalar computable from labeled training data alone, that reliably predicts whether prototype-based rescue will improve rare-class F1. When $S > 1.3$ and the scANVI baseline is miscalibrated ($\text{F1} < 0.75$), the three-stage pipeline---prototype scoring, geometric gate, and marker-gene verification---achieves consistent gains of $+23$ to $+78$ percentage points across immune, pancreatic, spleen, and liver datasets. When $S < 1.1$, the method correctly abstains, preserving the scANVI prediction with near-zero F1 change.

The most striking result is at the extreme of five labeled rare-class training cells: cDC1 F1 improves from 0.003 (scANVI) and 0.000 (kNN) to 0.986 (scRareRefine), demonstrating that latent-space geometry can recover rare-cell structure where both softmax and distance-voting classifiers fail. This data-efficiency advantage is most pronounced in the regime most relevant to discovery settings---where annotation budgets are tight and rare cells have not yet been systematically characterized.

\paragraph{Limitations and future work.}
The current evaluation is limited to nine rare cell types across seven datasets; broader benchmarking, including imbalance-aware retraining methods (scBalance, sc-SynO) and foundation model embeddings (scGPT), is needed to establish the generality of these findings. The separability threshold values (1.1 and 1.3) should be validated on independent datasets. Future work will also explore adapting the pipeline to the \emph{de novo} discovery setting, where the rare cell type is not specified in advance and the framework must jointly identify and characterize novel populations.

\paragraph{Data and code availability.}
All code, configuration files, and experiment scripts are available at [repository URL]. Raw data are available from the respective original studies.
